\section{A compact-family local principle and noncircular tables}
\label{sec:compact-family-local-principle}
We formulate the consequences of the action and finite-cover arguments
before specializing them. The hypotheses concern collision geometry,
not a local limit theorem for an induced map. The original circular
family is contained in the explicit family at the end of the section.

\subsection{Admissible compact families}
Let $\mathcal K$ be a compact parameter space. For $\lambda\in\mathcal K$
place translates of a strictly convex scatterer $O_\lambda$ at the
points of a fixed planar lattice $\Lambda$, with one scatterer per
fundamental cell. Use Euclidean billiard reflection and unit speed;
lattice coordinates, rather than physical Cartesian coordinates, are
used for displacement. Denote the perimeter by $L_\lambda$ and
normalize its arclength coordinate by $\alpha=2\pi r/L_\lambda$.
The collision probability is then the common measure
\[
 d\nu=(4\pi)^{-1}\cos\varphi\,d\alpha\,d\varphi.
\]
We impose the following geometric conditions.

\paragraph{Uniform table geometry.}
The boundary parametrizations form a continuous $C^3$ family with
positive lower and upper curvature bounds, a positive separation
between distinct scatterers, and a common finite horizon. The free
region in every finite lattice cover is connected. This last condition
is topological; it is not used as a synonym for mixing.

\paragraph{Local spaces and unweighted geometric estimates.}
On a finite parameter cover there are local common collision spaces
with the norms of Section~\ref{sec:action-norm-details}, faithful strong
completion, compact strong-to-weak embedding and the unweighted
spectral splitting. The stable contraction, distortion, normalized
Jacobian and length sums, and product absolute continuity are those
specified in that section. Constants are uniform on the finite cover;
one global Banach space over all parameters is not required.

\paragraph{Action and matching.}
The identity $d\tau=T^*\vartheta-\vartheta$ holds on regular branches,
and the matched connectors remain regular at every intermediate
iterate as in \eqref{eq:v35-action-telescope} and
\eqref{eq:v35-weight-match-detail}. The unweighted matching partition
contains the required singularity and homogeneity cuts.

\paragraph{Finite-cover dynamics.}
For each prime $p$, the collision map of the actual $p\Lambda$ cover
is mixing for its normalized collision measure. This is an independent
hypothesis in the general principle. No mixing rate uniform in $p$
is required. Proposition~\ref{prop:v37-finite-cover-mixing} verifies
it for the circular and oriented-ellipse applications by an explicit
correlation-decay theorem, not from connectedness alone.

\paragraph{Parameter continuity and normalization.}
The collision operators are strongly-to-weakly continuous in the
local spaces. The preceding-flight comparison below supplies that
continuity for the twists. The mean roof is continuous, bounded above,
and bounded away from zero. Finite-lag covariance integrands converge
almost everywhere and have common integrable bounds. Together with
uniform summable correlation tails, this is the covariance continuity
used in the compactness argument.

\paragraph{Preceding-flight regularity.}
The preceding displacement and roof satisfy the partition conditions
\eqref{eq:v35-partition-boundary}, with a common number of pieces and
a uniform piecewise $C^{1/2}$ norm. On matched one-step branches the
roof is jointly $1/2$-H\"older in the parameter and state in local
coordinates. Finite itineraries and these functions converge outside
the collision singularities as the parameter varies. No estimate on
the complexity of an arbitrary long forward occupation word is part
of this condition.

\paragraph{Regular observation geometry.}
Fix a convex polygonal fundamental cell. The labels of the cells met
by a single flight belong to a fixed finite set. Write
$I_{\lambda,d}(x)\subset[0,\tau_\lambda(x))$ for the ages at which a
flight from collision state $x$ lies in cell $d$. The endpoints of
these intervals vary continuously, locally uniformly outside closed
neighborhoods of a finite union of null curves and boundary pieces.
For every convergent parameter sequence the neighborhoods can be
chosen with arbitrarily small limiting collision measure and with
null boundary. Include grazing, singular flights, cell tangencies and
vertex crossings among the exceptional sets. Empty intervals are
allowed. This is a condition on one flight, not on the long orbit.

Set $f^{\rm c}_\lambda=(\kappa_\lambda,
\tau_\lambda-\bar\tau_\lambda)$, where
$\bar\tau_\lambda=\int\tau_\lambda d\nu$. Time reversal gives
$\int\kappa_\lambda d\nu=0$. Let $\Sigma_\lambda$ be its collision
Green--Kubo matrix and put
\begin{equation}\label{eq:v35-family-covariances}
 D^{\rm clk}_\lambda=\operatorname{diag}(1,1,-1/\bar\tau_\lambda),
 \qquad
 \mathcal V_\lambda=\bar\tau_\lambda^{-1}
       D^{\rm clk}_\lambda\Sigma_\lambda(D^{\rm clk}_\lambda)^{\mathsf T}.
\end{equation}
The matrix $D^{\rm clk}$ is a clock-change map, not the four-coordinate
return covariance $D_R$ used elsewhere in the article.

\begin{theorem}[Uniform local principle for an action family]
\label{thm:v35-action-family-LLT}
For a family satisfying these conditions, $\Sigma_\lambda$ and
$\mathcal V_\lambda$ exist, are continuous and uniformly positive
definite. Under the stationary suspension probabilities
$d\mu_\lambda(x,a)=\bar\tau_\lambda^{-1}d\nu(x)\,da$, for every
$M<\infty$,
\begin{equation}\label{eq:v35-family-singleton}
 \sup_{\substack{\lambda\in\mathcal K,\ k\in\Z^2,\ m\ge0\\
                 |\xi_{t,\lambda}(k,m)|\le M}}
 \left|t^{3/2}\mu_\lambda\{W_\lambda(t)=k,C_\lambda(t)=m\}
                 -g_{\mathcal V_\lambda}(\xi_{t,\lambda}(k,m))\right|
                       \longrightarrow0,
\end{equation}
where $\xi_{t,\lambda}(k,m)=(k/\sqrt t,
(m-t/\bar\tau_\lambda)/\sqrt t)$.
The same formula with bounded regular initial and terminal flight-state
selectors has Gaussian amplitude equal to the product of their
stationary means. Here regular means jointly continuous bounded
families, or bounded families continuous off the closed exceptional
neighborhoods specified above, with uniform convergence on their
complements. Indicator selectors must have null limiting boundaries.
A uniform positive selected denominator is asserted only for
nonnegative selectors whose product of stationary means is bounded
below by a positive number.
\end{theorem}
\begin{proof}
We give the entire logical route and the order of its constants.
The local unweighted spectral splitting and
Lemma~\ref{lem:v35-multiplier-detail} realize
\[
 Q_{\lambda,z}=M_{\exp(i z\cdot(f^{\rm c}_\lambda\circ T_\lambda^{-1}))}
                                      \mathcal L_\lambda.
\]
The covariance series is exponentially summable: replace both factors
by their preceding-flight versions using invariance, and pair the
unweighted splitting between the two bounded strong multipliers.
For example each centered scalar component $v$ satisfies
\[
 \int v(v\circ T^j)d\nu
 =\ell\bigl(M_{v\circ T^{-1}}\mathcal N^j
                         M_{v\circ T^{-1}}\nu\bigr).
\]
Finite-lag continuity follows by dominated convergence outside the
finitely many singularity preimages. A finite parameter cover makes
the correlation tails uniform. Corollary~\ref{cor:v35-general-covariance}
now gives continuous uniform ellipticity. This argument does not
require a uniform size for the product rectangles on all finite covers.

For a fixed $B$, Proposition~\ref{prop:v35-expanded-LY} gives local
uniform power boundedness and quasi-compactness of $Q_{\lambda,\omega}$.
The parameter comparison at the end of
Section~\ref{sec:action-norm-details} gives strong-to-weak continuity.
Lemma~\ref{lem:v35-peripheral-density} and
Theorem~\ref{thm:v35-rotation-free-rigidity} exclude every nonzero
peripheral frequency in $\T^2\times[-B,B]$. Near zero, analytic
perturbation of the simple eigenvalue and differentiation of its
transfer pairing give
\begin{equation}\label{eq:v35-family-eigenvalue}
 \log\lambda^{\rm sp}_\lambda(z)
       =-\tfrac12z^{\mathsf T}\Sigma_\lambda z+O(|z|^3),
                  \qquad \Pi_\lambda(0)=\nu\ell.
\end{equation}
The linear term vanishes by centering, and the second derivative is
exactly the covariance sum just established. Resolvent contours on
a finite parameter cover give uniform errors near zero. On any fixed
compact frequency set avoiding zero, the strong/weak perturbation
argument and a finite covering give a common spectral radius less
than one, as in Theorem~\ref{thm:compact-collision-spectrum}.

Here are the measure-theoretic consequences, with no growing-band
substitution. For a time test $q$ with integrable compactly supported
Fourier transform, torus orthogonality gives the endpoint pairing
\begin{align*}
 &\int a(x)d(T_\lambda^m x)\one_{\{K_m=k\}}q(S_m-t)d\nu(x)\\
 &=\frac1{(2\pi)^3}\int_{\T^2\times\R}
   e^{-i[u\cdot k+b(t-m\bar\tau_\lambda)]}\widehat q(-b)
                 \ell(M_dQ_{\lambda,(u,b)}^m(a\nu))\,du\,db.
\end{align*}
The minus sign in $\widehat q(-b)$ follows from the argument $S_m-t$.
For fixed support the complement of zero is exponentially small;
rescaling $\sqrt m(u,b)$ near zero and using
\eqref{eq:v35-family-eigenvalue} gives the Gaussian with endpoint
amplitude $\nu(a)\nu(d)\int q$.
For an interval $J$, positive upper and lower Fourier envelopes of
band $B$ have masses $|J|\pm2\pi/B$. With nonnegative smooth endpoints
multiply their pointwise inequalities by the nonnegative source
measure. First take $m\to\infty$ with $B$ fixed, then let $B\to\infty$.
This yields the unsmoothed interval limit; one fixed majorant also
bounds all local masses by $Cm^{-3/2}$. Approximation and decomposition
into real and imaginary parts extend the endpoint class to continuous
complex functions. This is the argument of
Theorem~\ref{thm:two-endpoint-collision-LLT}, now with uniform local
inputs on $\mathcal K$.

In particular, along $\lambda_j\to\lambda_0$ and
$(k_j,t_j-m_j\bar\tau_{\lambda_j})/\sqrt{m_j}\to\zeta$, the positive
endpoint measures converge vaguely on bounded time slabs to
\begin{equation}\label{eq:v35-family-endpoint-measure}
 m_j^{3/2}(x,S_{m_j}-t_j,T_{\lambda_j}^{m_j}x)_*
                (\one_{\{K_{m_j}=k_j\}}\nu)
          \ \rightharpoonup\ g_{\Sigma_{\lambda_0}}(\zeta)\nu\otimes ds\otimes\nu.
\end{equation}
For a varying overlap test use a closed exceptional neighborhood of
small limiting measure and null boundary. Portmanteau bounds the
source mass of this closed set; the local mass bound and uniform
convergence on the complement control the other part. This is why
small measure of an arbitrary open set, or pointwise convergence of
tests alone, is not the regularity condition.

The exact physical identity is a finite sum over the two cell offsets.
Its $(d,e)$ term uses the collision label $k+d-e$ and the nonnegative
test
\[
 F^{d,e}_\lambda(x,s,y)=
   \int\one_{I_{\lambda,d}(x)}(a)
             \one_{I_{\lambda,e}(y)}(a-s)\,da.
\]
Indeed $s=S_m-t$ and the final age is $a-s$. The original count event
is exactly $0\le a-s<\tau_\lambda(y)$, so no boundary flight is
replaced by a completed-return surrogate. The observation geometry
makes these bounded, compactly supported tests admissible in
\eqref{eq:v35-family-endpoint-measure}. Their integrated amplitude is
\[
 \sum_{d,e}\int F^{d,e}_\lambda\,d\nu(x)\,ds\,d\nu(y)
       =\left(\int\tau_\lambda d\nu\right)^2=\bar\tau_\lambda^2.
\]
The source normalization is $1/\bar\tau_\lambda$; hence on the collision
scale the main term is $\bar\tau_\lambda g_{\Sigma_\lambda}$.
The finite label offsets change the central coordinate by
$O(m^{-1/2})$. The physical clock conversion has
\[
 \det\mathcal V_\lambda=\bar\tau_\lambda^{-5}\det\Sigma_\lambda,
            \qquad m/t\longrightarrow\bar\tau_\lambda^{-1}.
\]
This converts the collision-scale expression exactly to
\eqref{eq:v35-family-singleton}. Sequential compactness in $\lambda$
and the central coordinates proves its stated uniformity.

For selectors insert $\chi_\lambda(x,a)\psi_\lambda(y,a-s)$ before
integrating over $a$. Fubini gives amplitude
$\bar\tau_\lambda^2\mu_\lambda(\chi_\lambda)
\mu_\lambda(\psi_\lambda)$. The same local-measure argument applies
under the stated regularity. Uniform ellipticity and the positive
product-of-means condition give the claimed selected lower bound.
\end{proof}

\subsection{An explicit three-parameter noncircular family}
Keep the triangular lattice of the original problem and set
\begin{equation}\label{eq:v35-ellipses}
 O_{a,b,\theta}=\operatorname{Rot}_\theta
       \operatorname{diag}(a,b)\{z:|z|<1\},\qquad
 (a,b,\theta)\in[9/20,47/100]^2\times(\R/2\pi\Z).
\end{equation}
The case $a=b=R$ is exactly the original disk table; $a\ne b$ in
general removes its sixfold rotational symmetry. Rotation here is a
parameter of the obstacle, not an affine conjugacy of reflection.

\begin{lemma}[Uniform checks for the ellipse family]
\label{lem:v35-ellipse-geometry}
The family \eqref{eq:v35-ellipses} satisfies all the separated
geometric and dynamical conditions above. Its mean roof is
\begin{equation}\label{eq:v35-ellipse-mean}
 \bar\tau_{a,b,\theta}
       =\frac{\pi(\sqrt3/2-\pi ab)}{L_{a,b}},\qquad
 2\pi(9/20)\le L_{a,b}\le2\pi(47/100).
\end{equation}
All collision probabilities are the same in normalized arclength
coordinates.
\end{lemma}
\begin{proof}
Every ellipse contains the disk of radius $9/20$ and is contained in
the disk of radius $47/100$. Different obstacles are therefore
separated by at least $3/50$. The smallest circular table has finite
horizon, and containment transfers its upper flight bound to every
ellipse. The curvature in the usual ellipse parameter $v$ is
\[
 \frac{ab}{(a^2\sin^2v+b^2\cos^2v)^{3/2}}.
\]
It lies between $(9/20)^2/(47/100)^3$ and
$(47/100)^2/(9/20)^3$. Its boundary derivatives and normalized
arclength changes have common $C^3$ bounds, and the analytic family
also has the uniform fourth-derivative bounds required by \cite{SYZ}.
The positive gaps give a connected free region in every finite cover.
Mixing is a separate consequence of
Proposition~\ref{prop:v37-finite-cover-mixing}, with its precise
collision-map correlation theorem and transverse-horizon check; it is
not deduced from connectedness. The local product construction is the
unweighted one in \cite[Section 8]{Young}. A rectangular Euclidean
cover may be used without an affine change of the reflection law.
The mean-roof
formula is the collision suspension identity
$\bar\tau=\pi\,\operatorname{area}(\text{free cell})/L$.
The area is $\sqrt3/2-\pi ab$; the perimeter bounds follow from the
speed of the ellipse parametrization. This proves
\eqref{eq:v35-ellipse-mean} and the common probability normalization.

We verify the preceding-flight condition directly. Let $z$ be the
current boundary point, $w$ the unit backwards incoming velocity and
$d$ a candidate preceding center. Put
$A=\operatorname{Rot}_\theta\operatorname{diag}(a^{-2},b^{-2})
\operatorname{Rot}_\theta^{\mathsf T}$. The candidate entry length is
\begin{equation}\label{eq:v35-ellipse-entry-root}
 \ell_d=\frac{-B_0-\sqrt{B_0^2-A_0C_0}}{A_0},\quad
 A_0=w^{\mathsf T}Aw,\quad B_0=w^{\mathsf T}A(z-d),\quad
 C_0=(z-d)^{\mathsf T}A(z-d)-1.
\end{equation}
For an obstacle other than the current one, $C_0>0$ and $A_0>0$.
A positive entry requires $B_0<0$ and $B_0^2-A_0C_0\ge0$.
The displayed minus-square-root solution is the first intersection;
the plus solution is the exit, and is not selected as the entry.
For the current obstacle $C_0=0$ and the backward outgoing ray has
$B_0>0$ off grazing, so its roots are zero and a negative value.
The zero root is discarded. Finite horizon
bounds all possible $d$ in one fixed finite set. The coefficients are
smooth and uniformly bounded in state and parameter, with
$A_0\ge(47/100)^{-2}$. Thus the entry root is uniformly $1/2$-H\"older
up to tangency, by the square-root inequality. Its preceding lattice
label is $-d$.

For clarity, finiteness of candidates by itself is not the multiplier
partition check. In rational trigonometric charts for the boundary,
velocity and orientation, the entry inequalities, tangency equation
and comparisons of the finitely many roots are semialgebraic of fixed
complexity. Uniform cell decomposition splits their boundaries into
a bounded number of monotone arcs, uniformly in $(a,b,\theta)$.
Distinct first-hit regions can meet only on a backward singularity:
a tie of two regular first entries would hit two separated obstacles
at the same point. Backward singularity arcs are unstable-cone arcs,
with common transversality to stable curves. A stable graph consequently
crosses each such monotone arc at most once after the finite chart
refinement. The length of its part within distance $\epsilon$ of an
arc is at most $C\epsilon$, by the slope separation. The coordinate
cuts have the same property. Intersecting this fixed finite partition
with any homogeneity strip preserves a common component and
intersection bound. Normalized arclength is a uniformly smooth
monotone coordinate change and preserves these properties. This proves
both \eqref{eq:v35-partition-boundary} and the uniform partition count,
not merely piecewise continuity of the root.

Finally, the bounded straight segment of a single flight intersects
each translated convex cell in an interval, and meets only a fixed
finite number of cells. Away from a tangent or a vertex crossing its
entry and exit ages are continuous functions of its endpoints and
direction. Changes of collision root, grazing, vertex crossings and
flights lying in a cell edge are finite unions of semialgebraic curves
or boundary pieces in collision coordinates. Each is null for $\nu$;
for a flight along a cell edge the fixed direction condition already
has zero two-dimensional measure. On compact complements of their
closed neighborhoods all interval endpoints vary uniformly. Choose
neighborhood radii at continuity points of their measure; their
boundaries are then null and their measures tend to zero. Along any
converging parameter sequence the corresponding regular roots persist
on this compact complement. This gives the stated varying-test
condition, including the appearance and disappearance of empty cell
intervals. The same construction handles the products of the two
endpoint spaces and a bounded age slab.
\end{proof}

\begin{corollary}[The microscopic law for oriented ellipses]
\label{cor:v35-ellipse-local-law}
The full conclusions of Theorem~\ref{thm:v35-action-family-LLT} hold
uniformly on \eqref{eq:v35-ellipses}. In particular they hold along
arbitrary sequences of eccentricities and orientations converging to
a circle, as well as within a noncircular family. The original circular
stationary theorem is its restriction to $a=b$.
\end{corollary}
\begin{proof}
Apply Lemma~\ref{lem:v35-ellipse-geometry} and the theorem. In
particular finite-cover arithmetic, rather than a rotation of a circle,
is what proves nondegeneracy and peripheral exclusion here.
\end{proof}

\subsection{Conditioning on the unchanged trajectory}
\begin{corollary}[Uniform posterior approximation in the same family]
\label{cor:v35-family-posterior}
On a fixed central compact set the microscopic event in
\eqref{eq:v35-family-singleton} has probability at least $c_Mt^{-3/2}$
for all sufficiently large $t$, uniformly in the parameter.
For the positive finite-band likelihood of
Theorem~\ref{thm:physical-time-selector-reconstruction}, formed with the same
physical flight-age indicator, the normalized source total-variation
error is at most $C_M m^{3/2}/B$. In particular $B=m^{P+3/2}$ gives
$O(m^{-P})$ for every fixed $P>0$, against arbitrary bounded measurable
tests on the original trajectory space. This assertion includes the
ellipse family and involves no growing-band spectral estimate.
\end{corollary}
\begin{proof}
Continuity and uniform ellipticity bound the Gaussian away from zero
on the compact set; the local theorem gives the denominator. The
positive smoothing kernel used in
Theorem~\ref{thm:physical-time-selector-reconstruction} has a finite first
moment of order $B^{-1}$. Translation by $v$ changes the indicator of
one age interval in $L^1(da)$ by at most $2|v|$. There are only uniformly
many initial and terminal cell intervals, and $\bar\tau$ is bounded
below. The identical source calculation therefore gives unnormalized
error at most $C/B$ throughout the present family. If finite positive
measures $\alpha,\beta$ have masses $a,b>0$, then
$\|\alpha/a-\beta/b\|_{\rm TV}\le2\|\alpha-\beta\|_{\rm TV}/a$,
with total variation taken as the norm of a signed measure.
Use the proved microscopic mass and $t\asymp m$. The local theorem
was obtained with a fixed band first; the later choice of $B$ concerns
only the explicit smoothing error. It provides a posterior comparison,
not a Gaussian amplitude for an arbitrary path selector.
\end{proof}
